% Qualitative example 1 (SIR-4 Biology, cross-field): protein-RNA binding-site query whose gold
% inspiration is a computer-vision paper. Numbers from colab_qualitative_biology.ipynb (8 Sep 2026):
%   rank of the gold inside SciGraphIR: raw Qwen3 cosine 113 -> multi-view scorer 13 -> graph reasoner 3 -> fused 8
%   baselines: BM25 79, BGE-large 117, Qwen3-Embedding 113, SPECTER2 >300, SciNCL >300, ReasonIR-8B 70
%   OpenIE entity-graph model: graph channel 3329 (fused 13)
%
% Preamble needs:  \usepackage{xcolor}  \usepackage[most]{tcolorbox}  \usepackage{booktabs}
% Colours are defined here with \providecolor so this file can be \input anywhere.

\providecolor{hlblue}{RGB}{214,228,247}
\providecolor{hlgreen}{RGB}{214,240,224}
\providecolor{hlorange}{RGB}{252,232,200}
\providecolor{hlpurple}{RGB}{228,220,246}
\providecommand{\hl}[2]{\colorbox{#1}{\strut #2}}
\providecommand{\qbox}[2]{%
  \begin{tcolorbox}[enhanced, colback=gray!4, colframe=black!70, boxrule=0.6pt, arc=2.5pt,
                    left=6pt, right=6pt, top=4pt, bottom=4pt, before skip=4pt, after skip=4pt]
  \textbf{#1}\quad #2
  \end{tcolorbox}}

\begin{figure}[t]
\centering
\small
\qbox{Research problem}{``How can a computational model for \hl{hlblue}{protein--RNA binding site prediction}
integrate global and local features, \hl{hlorange}{handle class imbalance}, and produce
\hl{hlgreen}{well-calibrated predictions}? Existing methods extract features at a single scale, and because
binding-site data are \hl{hlorange}{inherently imbalanced} they trade precision against recall and their
confidence scores are poorly calibrated.''}

\qbox{Hypothetical answer $h_1$}{``An \hl{hlpurple}{ensemble of classifiers trained on balanced subsets} of the
data, using \hl{hlorange}{SMOTE to mitigate class imbalance} and improve prediction accuracy.''\quad
{\footnotesize (best-matching view: cosine to the gold 0.62, against 0.44 for the raw query)}}

\qbox{Retrieved inspiration}{\emph{An Optimal Transport View of Class-Imbalanced Visual Recognition}
(computer vision).\quad
{\footnotesize rank of this paper: Qwen3 cosine \textbf{113} $\rightarrow$ multi-view scorer \textbf{13}
$\rightarrow$ graph reasoner \textbf{3} $\rightarrow$ fused \textbf{8}; \; BGE-large 117, SPECTER2 $>$300,
ReasonIR-8B 70.}}

\caption{A cross-field retrieval on SIR-4 Biology. The query never names optimal transport; the hypothetical
answer written for it names the \emph{function} the query needs (\hl{hlorange}{handling class imbalance}),
which the scorer matches against the vision paper, and the graph reasoner then reaches the paper through the
limitation frame shared by the query and the paper (Table~\ref{tab:path-interpretations}).}
\label{fig:qual-example-ot}
\end{figure}

\begin{table}[t]
\centering
\resulttablestyle
\setlength{\tabcolsep}{5pt}
\renewcommand{\arraystretch}{1.25}
\begin{tabular}{@{}p{0.13\linewidth} p{0.85\linewidth}@{}}
\toprule
\textbf{Query} & How can a computational model for protein--RNA binding site prediction integrate global and
local features, handle class imbalance, and produce well-calibrated predictions? \\
\textbf{Gold} & \emph{An Optimal Transport View of Class-Imbalanced Visual Recognition} \\
\textbf{Ranks} & Qwen3 cosine 113 \;|\; multi-view scorer 13 \;|\; graph reasoner 3 \;|\; \textsc{SciGraphIR} 8
\;|\; BGE-large 117 \;|\; ReasonIR-8B 70 \\
\midrule
\textbf{Paths} \newline (SciGraph frame graph) &
12.44:\; (\textsc{limitation}: sensitivity to class imbalance,\; \textit{overcomes}$^{-1}$,\;
\textsc{method}: optimal transport-based method) $\rightarrow$
(\textsc{method}: optimal transport-based method,\; \textit{contributes}$^{-1}$,\; \textbf{gold paper}) \newline
11.31:\; (\textsc{function}: improve classification on imbalanced data,\; \textit{achieves}$^{-1}$,\;
\textsc{method}: optimal transport-based method) $\rightarrow$
(\textsc{method}: optimal transport-based method,\; \textit{contributes}$^{-1}$,\; \textbf{gold paper}) \\
\addlinespace
\textbf{Path} \newline (OpenIE entity graph) &
0.21:\; (deep learning,\; \textit{is used in},\; computer vision) $\rightarrow$
(computer vision,\; \textit{is applied in}$^{-1}$,\; optimal mass transport) $\rightarrow$
(optimal mass transport,\; \textit{equivalent}$^{-1}$,\; optimal transport) $\rightarrow$
(optimal transport,\; \textit{is mentioned in},\; \textbf{gold paper}) \newline
{\footnotesize graph channel rank of the gold under the OpenIE model: 3{,}329} \\
\bottomrule
\end{tabular}
\caption{Path interpretations for the query of Figure~\ref{fig:qual-example-ot}. Paths are the top-weighted
routes from the query's seed frames to the gold under the trained reasoner, found by beam search over the
gradient of the graph score with respect to each layer's edge weights (NBFNet, GFM-RAG); $r^{-1}$ is the inverse
relation. On the frame graph the query's own \textsc{limitation} and \textsc{function} frames link in two hops to the
method the gold paper contributes. On the OpenIE entity graph built from the same corpus the only route is a chain
of generic co-mentions, and the graph channel ranks the paper at 3{,}329.}
\label{tab:path-interpretations}
\end{table}
